% Part: first-order-logic
% Chapter: models-of-arithmetic
% Section: models-of-pa

\documentclass[../../../include/open-logic-section]{subfiles}

\begin{document}

\olfileid[pt]{mod}{mar}{mpa}
\section{Modelos de $\Th{PA}$}

\begin{explain}
Todo modelo não padrão de~$\Th{TA}$ também é um modelo de~$\Th{PA}$.
Sabemos que existem modelos não padrão de~$\Th{TA}$ e, portanto,
de~$\Th{PA}$. Sabemos ainda que esses modelos não padrão contêm
``números'' não padrão, isto é, !!{element}s do domínio que estão
``além'' de todos os ``números'' padrão. Mas como eles se organizam?
Quantos existem? Vimos que modelos da teoria mais fraca~$\Th{Q}$ podem
conter apenas um único número não padrão. Essas !!{structure}s simples,
porém, não são modelos de $\Th{PA}$ nem de $\Th{TA}$.

A chave para compreender a estrutura dos modelos de $\Th{PA}$ ou
$\Th{TA}$ é examinar quais fatos são !!{derivable} nessas teorias. Por
exemplo, $\Th{PA}$ já demonstra $\lforall[x][\eq/[x][x']]$ e
$\lforall[x][\lforall[y][\eq[(x+y)][(y+x)]]]$; isso exclui como modelos
de~$\Th{PA}$ as estruturas simples nas quais essas !!{sentence}s são
falsas.

Suponha que~$\Struct{M}$ seja um modelo de~$\Th{PA}$. Então, se
$\Th{PA} \Proves !A$, temos $\Sat{M}{!A}$. Voltemos a usar $\nszero$
para $\Assign{\Obj{0}}{M}$, $\nssucc$ para $\Assign{\prime}{M}$,
$\nsplus$ para $\Assign{+}{M}$, $\nstimes$ para
$\Assign{\times}{M}$ e $\nsless$ para $\Assign{<}{M}$. Toda
!!{sentence}~$!A$ enuncia, assim, alguma condição sobre $\nszero$,
$\nssucc$, $\nsplus$, $\nstimes$ e $\nsless$; se $\Sat{M}{!A}$, essa
condição deve ser satisfeita. Por exemplo, se $\Sat{M}{!Q_1}$, isto é,
$\Sat{M}{\lforall[x][\lforall[y][(\eq[x'][y'] \lif \eq[x][y])]]}$,
então $\nssucc$ deve ser !!{injective}.
\end{explain}

\begin{prop}
Em $\Struct{M}$, $\nsless$ é uma ordem linear estrita, isto é,
satisfaz:
\begin{enumerate}
\item Não vale $x \nsless x$ para nenhum~$x \in \Domain{M}$.
\item Se $x \nsless y$ e $y \nsless z$, então $x \nsless z$.
\item Para quaisquer $x \neq y$, vale $x \nsless y$ ou $y \nsless x$.
\end{enumerate}
\end{prop}

\begin{proof}
$\Th{PA}$ demonstra:
\begin{enumerate}
\item $\lforall[x][\lnot x < x]$
\item $\lforall[x][\lforall[y][\lforall[z][((x < y \land y < z) \lif x < z)]]]$
\item $\lforall[x][\lforall[y][((x < y \lor y < x) \lor \eq[x][y]))]]$
\end{enumerate}
\end{proof}

\begin{prop}
\ollabel{prop:M-discrete} $\nszero$ é o menor !!{element}
de~$\Domain{M}$ na ordenação $\nsless$. Para todo $x$, $x \nsless
x^\nssucc$, e $x^\nssucc$ é o menor !!{element} segundo $\nsless$ com
essa propriedade. Para todo $x$, existe um único $y$ tal que
$y^\nssucc = x$. (Chamamos $y$ de ``predecessor'' de~$x$ em
$\Struct{M}$ e o denotamos por~$^\nssucc x$.)
\end{prop}

\begin{proof}
Exercício.  
\end{proof}

\begin{prob}
Encontre !!{sentence}s em~$\Lang{L_A}$ !!{derivable}s em~$\Th{PA}$ (e,
portanto, verdadeiras em~$\Struct{N}$) que garantam as propriedades de
$\nszero$, $\nssucc$ e $\nsless$ em
\olref[mod][mar][mpa]{prop:M-discrete}
\end{prob}

\begin{prop}
Todos os !!{element}s padrão de~$\Struct{M}$ são menores (segundo
$\nsless$) que todos os !!{element}s não padrão.
\end{prop}

\begin{proof}
Usaremos $n$ como abreviação de $\Value{\num{n}}{M}$, um
!!{element} padrão de~$\Struct{M}$. $\Th{Q}$ já demonstra que, para
todo~$n \in \Nat$, $\lforall[x][(x < \num{n}' \lif
(\eq[x][\num{0}] \lor \eq[x][\num{1}] \lor \dots \lor
\eq[x][\num{n}]))]$. Não há !!{element}s que sejam $\nsless
\nszero$. Portanto, se $n$ é padrão e $x$ é não padrão, não podemos
ter $x \nsless n$. Por definição, um elemento não padrão é aquele que
não é $\Value{\num{n}}{M}$ para nenhum~$n \in \Nat$, logo também
$x \neq n$. Como $\nsless$ é uma ordem linear, devemos ter
$n \nsless x$.
\end{proof}

\begin{prop}
Todo !!{element} não padrão~$x$ de~$\Domain{M}$ pertence ao subconjunto
\[
\dots ^{\nssucc\nssucc\nssucc}x \nsless ^{\nssucc\nssucc}x \nsless
^{\nssucc}x \nsless x \nsless x^{\nssucc} \nsless x^{\nssucc\nssucc}
\nsless x^{\nssucc\nssucc\nssucc} \nsless \dots
\]
Chamamos esse subconjunto de \emph{bloco de~$x$} e o escrevemos como
$[x]$. Ele não tem menor nem maior !!{element}. Pode ser caracterizado
como o conjunto dos $y \in \Domain{M}$ tais que, para algum~$n$ padrão,
$x \nsplus n = y$ ou $y \nsplus n = x$.
\end{prop}

\begin{proof}
É claro que um conjunto~$[x]$ assim sempre existe, pois todo
!!{element}~$y$ de $\Domain{M}$ tem um único sucessor~$y^\nssucc$ e um
único predecessor~$^\nssucc y$. Para !!{element}s sucessivos $y$,
$y^\nssucc$, temos $y \nsless y^\nssucc$, e $y^\nssucc$ é o
!!{element} $\nsless$-mínimo de~$\Domain{M}$ tal que $y$ é
$\nsless$-menor que ele. Como
sempre $^\nssucc y \nsless y$ e $y \nsless y^\nssucc$, $[x]$ não tem
menor nem maior !!{element}. Se $y \in [x]$, então $x \in [y]$, pois
nesse caso ou $y^{\nssucc\dots\nssucc} = x$ ou
$x^{\nssucc\dots\nssucc} = y$. Se $y^{\nssucc\dots\nssucc} = x$ (com
$n$ ocorrências de $\nssucc$), então $y \nsplus n = x$, e
reciprocamente, pois $\Th{PA} \Proves
\lforall[x][\eq[x^{\prime\dots\prime}][(x + \num{n})]]$ (se $n$ é o
número de ocorrências de~$\prime$).
\end{proof}

\begin{prop}
Se $[x] \neq [y]$ e $x \nsless y$, então, para quaisquer $u \in [x]$
e $v \in [y]$, temos $u \nsless v$.
\end{prop}

\begin{proof}
Observe que $\Th{PA} \Proves \lforall[x][\lforall[y][(x < y \lif (x' < y
    \lor x' = y))]]$. Portanto, se $u \nsless v$, também temos
$u \nsplus n^\nssucc \nsless v$ para todo~$n$, desde que
$[u] \neq [v]$.

Todo $u \in [x]$ é $\nsless y$: por hipótese, $x \nsless y$.
Se $u \nsless x$, então $u \nsless y$ por transitividade. E, se
$x \nsless u$, mas $u \in [x]$, temos $u = x\nsplus n^\nssucc$ para
algum~$n$, logo $u \nsless y$ pelo fato que acabamos de demonstrar.
  
Suponha agora que $v \in [y]$ seja $\nsless y$, isto é,
$v \nsplus m^\nssucc = y$ para algum~$m$ padrão. Isso exclui
$v \nsless x$; caso contrário, $y = v \nsplus m^\nssucc \nsless x$.
Também é claro que $x \neq v$; caso contrário,
$x \nsplus m^\nssucc = v \nsplus m^\nssucc = y$, e teríamos
$[x] = [y]$. Logo, $x \nsless v$. Mas então também
$x \nsplus n^\nssucc \nsless v$ para todo~$n$. Portanto, se
$x \nsless u$ e $u \in [x]$, temos $u \nsless v$. Se $u \nsless x$,
então $u \nsless v$ por transitividade.

Por fim, se $y \nsless v$, então $u \nsless v$, pois, como mostramos,
$u \nsless y$ e $y \nsless v$.
\end{proof}

\begin{cor}
Se $[x] \neq [y]$, então $[x] \cap [y] = \emptyset$.
\end{cor}

\begin{proof}
Suponha que $z \in [x]$ e $x \nsless y$. Então $z \nsless u$ para
todo $u \in [y]$. Se $z \in [y]$, teríamos $z \nsless z$. O caso em
que $y \nsless x$ é análogo.
\end{proof}

\begin{explain}
Isso significa que os próprios blocos podem ser ordenados de uma
maneira que respeite $\nsless$: $[x] \nsless [y]$ se, e somente se,
$x \nsless y$, ou, de modo equivalente, se $u \nsless v$ para
quaisquer $u \in [x]$ e $v \in [y]$. É claro que o bloco padrão $[0]$
é o menor bloco. Ele não intersecta nenhum bloco não padrão, e dois
blocos não padrão tampouco se intersectam. Em particular, não se pode
``alcançar'' outro bloco tomando repetidamente sucessores ou
predecessores.
\end{explain}

\begin{prop}
Se $x$ e $y$ são não padrão, então $x \nsless x \nsplus y$ e
$x \nsplus y \notin [x]$.
\end{prop}

\begin{proof}
Se $y$ é não padrão, então $y \neq \nszero$. $\Th{PA} \Proves
\lforall[x][(y \neq \Obj{0} \lif x < (x+y))]$. Suponha agora que
$x \nsplus y \in [x]$. Como $x \nsless x \nsplus y$, teríamos
$x \nsplus n^\nssucc = x \nsplus y$. Mas $\Th{PA} \Proves
\lforall[x][\lforall[y][\lforall[z][(\eq[(x+y)][(x+z)] \lif y = z)]]]$
(a lei do cancelamento para a adição). Isso significaria que
$y = n^\nssucc$ para algum~$n$ padrão; mas supomos que $y$ seja não
padrão.
\end{proof}

\begin{prop}
Não existe menor bloco não padrão.
\end{prop}

\begin{proof}
$\Th{PA} \Proves \lforall[x][\lexists[y][(\eq[(y+y)][x] \lor
      \eq[(y+y)'][x])]]$, isto é, todo $x$ é divisível por~$2$
  (possivelmente com resto~$1$). Se $x$ é não padrão, $y$ também é.
  Pela proposição anterior, $y \nsless y \nsplus y$ e $y \nsplus y
  \notin [y]$. Assim, também $y \nsless (y \nsplus y)^\nssucc$ e $(y
  \nsplus y)^\nssucc \notin [y]$. Mas $x = y \nsplus y$ ou $x = (y
  \nsplus y)^\nssucc$, logo $y \nsless x$ e $y \notin [x]$.
\end{proof}

\begin{prop}
Não existe maior bloco.
\end{prop}

\begin{proof}
Exercício.
\end{proof}

\begin{prob}
Mostre que, em um modelo não padrão de~$\Th{PA}$, não existe maior
bloco.
\end{prob}

\begin{prop}
\ollabel{prop:blocks-dense}
A ordenação dos blocos é densa. Isto é, se $x \nsless y$ e
$[x] \neq [y]$, então existe entre eles um bloco $[z]$ distinto de
ambos.
\end{prop}

\begin{proof}
Suponha que $x \nsless y$. Como antes, $x \nsplus y$ é divisível por
dois (possivelmente com resto): existe $z \in \Domain{M}$ tal que ou
$x \oplus y = z \oplus z$, ou $x \oplus y = (z \oplus z)^\nssucc$.
O elemento $z$ é a ``média'' de $x$ e~$y$, e $x \nsless z$ e
$z \nsless y$.
\end{proof}

\begin{prob}
Escreva uma demonstração detalhada de
\olref[mod][mar][mpa]{prop:blocks-dense}. Qual !!{sentence} $\Th{PA}$
deve !!{derive} para garantir a existência de~$z$? Por que
$x \nsless z$ e $z \nsless y$, e por que $[x] \neq [z]$ e
$[z] \neq [y]$?
\end{prob}

\begin{explain}
Os blocos não padrão são, portanto, ordenados como os racionais: eles
formam uma ordenação linear densa !!{denumerable} sem extremidades. Pode-se mostrar
que quaisquer duas ordenações !!{denumerable}s desse tipo são isomorfas. Segue-se
que, para quaisquer dois modelos não padrão !!{enumerable}s $\Struct{M}_1$ e
$\Struct{M_2}$ da aritmética verdadeira, seus redutos à linguagem que contém
apenas $<$ e $=$ são isomorfos. Com efeito, um isomorfismo $h$ pode ser
definido do modo seguinte: as partes padrão de $\Struct{M_1}$ e
$\Struct{M_2}$ são isomorfas ao modelo padrão $\Struct{N}$ e, portanto, entre
si. Os blocos que compõem a parte não padrão são eles próprios ordenados como
os racionais e, por conseguinte, isomorfos; um isomorfismo entre os blocos
pode ser estendido a um isomorfismo \emph{dentro} dos blocos emparelhando-se
elementos arbitrários em cada um e fazendo-se então corresponder a imagem do
sucessor de $x$ em $\Struct{M_1}$ ao sucessor da imagem de $x$ em
$\Struct{M_2}$. Observe que disso \emph{não} se segue que
$\mathfrak{M}_1$ e $\mathfrak{M}_2$ sejam isomorfos na linguagem completa da
aritmética (com efeito, o isomorfismo é sempre relativo a !!a{language}), pois
há maneiras não isomorfas de definir a adição e a multiplicação sobre
$\Domain{M_1}$ e $\Domain{M_2}$. (Isso também decorre de um célebre teorema
de Vaught segundo o qual o número de modelos enumeráveis de uma teoria
completa não pode ser~2.)
\end{explain}
\end{document}
